\chapter{Input Parsing and Validation}
\label{ch:validation}

\section{Netlist and Cell Library Grammar}
\label{sec:val-grammar}

The netlist format is a flat, line-oriented text grammar with three
statement forms, one per line, blank lines and \code{\#}-prefixed comments
ignored:

\begin{lstlisting}[language=bash, style=pythonstyle, numbers=none]
INPUT  <signal_name>
OUTPUT <signal_name>
<instance_name> <cell_name> <input_net_1> ... <input_net_N> <output_net>
\end{lstlisting}

The standard-cell library is a single JSON file, shared across every
circuit, describing thirty-six cells across twelve logic families
(\code{INV}, \code{BUF}, \code{AND2}, \code{NAND2}, \code{OR2}, \code{NOR2},
\code{XOR2}, \code{XNOR2}, three-input variants, and so on), each available
in three drive-strength sizes (\code{X1}, \code{X2}, \code{X4}). Every cell
entry supplies, per input pin, its capacitance, logical effort, and timing
sense (\cref{sec:bg-unate}), and, per cell, its intrinsic rise/fall delay,
rise/fall output resistance, internal capacitance, normalized parasitic
delay, leakage power, and area --- i.e.\ every quantity that appears on the
right-hand side of \cref{eq:cload,eq:drise,eq:dfall,eq:le-stage,eq:power-area,eq:power-total}.

\section{Configuration Schema}
\label{sec:val-config}

The configuration file supplies every circuit- and run-specific parameter
that is not part of the netlist or the (shared) cell library: primary-input
arrival times, primary-output required times and external loads, operating
conditions (supply voltage, clock frequency, temperature, switching
activity factors), area and power budgets, optimizer settings (allowed
sizes, iteration limit, cost weights), and Monte Carlo settings (sample
count, variation strengths, random seed). Every field is validated for
presence, type, and range before any analysis is attempted; a missing or
malformed field is rejected with a specific error message rather than
silently substituted with a default, on the reasoning that a silently
wrong configuration value is more dangerous than a loud rejection.

\section{Structural Validation Stages}
\label{sec:val-stages}

Validation proceeds through four ordered stages, and the process is
aborted at the first stage that fails so that no analysis is ever attempted
on an input the tool cannot trust.

\begin{enumerate}
    \item \textbf{Configuration.} Schema and range validation as described
          above.
    \item \textbf{Cell library.} Schema validation of the shared library
          file (every cell has a complete, correctly typed set of the
          fields listed in \cref{sec:val-grammar}).
    \item \textbf{Netlist.} Parses the grammar and performs every
          netlist-level structural check listed in \cref{tab:error-classes}.
    \item \textbf{DAG.} Builds the circuit graph and topologically sorts it;
          failure at this stage means a combinational (feedback) loop was
          detected.
\end{enumerate}

\subsection{Combinational Loop Detection}
\label{sec:val-cycle}

A combinational loop cannot be resolved by a linear topological ordering:
Kahn's algorithm repeatedly removes vertices with zero remaining in-degree
and terminates having placed every vertex if and only if the graph is
acyclic. If any vertices remain unplaced when the algorithm's queue empties,
those vertices participate in (or exclusively depend on) a cycle, and the
netlist is rejected with the specific set of implicated gate instances
reported to the user.

\section{Error Taxonomy and the Invalid-Circuit Benchmark Set}
\label{sec:val-errors}

Six structural error classes are validated against, matching the six
invalid benchmark circuits supplied with the assignment
(\code{e01}--\code{e06}). \Cref{tab:error-classes} summarizes each class,
the pipeline stage at which it is detected, and the diagnostic message
format produced.

\begin{longtable}{@{}p{0.25\textwidth}p{0.10\textwidth}p{0.5\textwidth}@{}}
\toprule
\textbf{Error class} & \textbf{Stage} & \textbf{Detection} \\
\midrule
\endhead
\code{undefined\_cell} &
  Netlist &
  A gate instance references a cell name absent from the library. \\[4pt]
\code{undriven\_signal} &
  Netlist &
  A gate input net is neither a primary input nor the output of any gate
  in the netlist. \\[4pt]
\code{multiple\_drivers} &
  Netlist &
  A net name appears as the output of more than one gate instance (or of
  both a gate and a declared primary input). \\[4pt]
\code{invalid\_input\_count} &
  Netlist &
  The number of input nets listed for a gate instance does not match the
  defined cell's number of input pins. \\[4pt]
\code{undriven\_output} &
  Netlist &
  A signal declared with \code{OUTPUT} is not the output of any gate
  instance. \\[4pt]
\code{combinational\_loop} &
  DAG &
  Topological sort fails to place every gate (\cref{sec:val-cycle}). \\
\bottomrule
\caption{Structural error taxonomy.}
\label{tab:error-classes}
\end{longtable}

